\section{引言}

2-bit大语言模型量化把连续权重压入稀疏离散空间。向量量化以一个码字表示多个权重，在相同比特预算下获得更强表达力，也把优化变量变成沿残差流、注意力和MLP非线性相互作用的assignment集合。AQLM、QuIP\#和QTIP说明表示几何、二阶敏感度与可实现编码同等重要\cite{aqlm,quip,qtip}；但部署对象最终是硬整数索引，soft或独立局部代理不会自动等价于硬模型行为。

本文从block-scaled shared-codebook VQ出发。Scale先把不同动态范围的block映射到共同领域，使整层共享4096个六维码字；Vector-GPTQ在top-128几何候选中加入激活二阶信息与误差反馈，形成硬锚点。锚点之后，当前码字8-way附近领域内的assignment是自然剩余自由度，可由binary Concrete概率训练，但真正问题是哪些硬切换值得被激活。

此前聚合一阶在98.20\%向量上找到下降邻居，四视图共识将资格率降到48.71\%，局部曲率又稳定优于共识；三者形成的硬集合仍全面输给source。V13不再调整单点分数，而把每个Linear固定top-128 bundle作为离散坐标，按layers 28--31因果顺序，在已接受状态下真实评估每层七个坐标。

四个条件坐标使训练目标从0.3298单调降至0.3208，held-out任务loss也改善1.09\%；然而文本CE恶化0.919\%并越过保护线。贡献为：（1）提出直接测量$F(S\cup B)-F(S)$的因果集合条件搜索，并在当前受限设置中缓解独立代理的交互计分缺口；（2）建立显式causal mask、同batch parity与曲率--block MSE正控制的可审计重放；（3）把机制边界从独立代理失真推进到跨域目标错位：task-only真实边际收益仍非语言保持证书。
